% ---------------------------------------------------------------------
%  CHƯƠNG 2 — CƠ SỞ LÝ THUYẾT VÀ CÔNG NGHỆ
% ---------------------------------------------------------------------
\chapter{Cơ sở lý thuyết và công nghệ}
\label{ch:coso}

\begin{muctieu}
Trình bày các kiến thức lý thuyết nền tảng gồm các tính năng mới của Java 21,
lập trình mạng Socket TCP, kỹ thuật đồng bộ hóa đa luồng, cơ chế phòng chống
Deadlock, cùng danh mục công cụ môi trường được sử dụng.
\end{muctieu}

\section{Kiến thức nền tảng Java 21 và Mạng TCP}

\begin{dinhnghia}[Giao thức văn bản hướng dòng (Stream-based Text Protocol)]
Là giao thức truyền thông dựa trên luồng ký tự thuần túy, trong đó các thông
điệp yêu cầu và phản hồi được phân tách bởi ký tự xuống dòng
(\code{\textbackslash n} hoặc \code{\textbackslash r\textbackslash n}). Mỗi
dòng chứa một từ khóa lệnh đứng đầu, theo sau là các tham số ngăn cách bằng
khoảng trắng.
\end{dinhnghia}

Ưu điểm vượt trội của giao thức văn bản là tính trong sáng, độc lập nền tảng,
dễ dàng kiểm thử thủ công bằng các công cụ như Telnet/Netcat và không phụ
thuộc vào phiên bản lớp như cơ chế tuần tự hóa nhị phân
(\code{Serializable})~\cite{bloch2018effective}.

\begin{dinhnghia}[Tranh chấp tài nguyên và Bế tắc luồng (Race Condition \& Deadlock)]
\textbf{Race Condition} xuất hiện khi hai hay nhiều luồng cùng truy cập và sửa
đổi dữ liệu dùng chung (ví dụ cùng đoán một số bí mật hoặc cùng cập nhật bảng
điểm) mà kết quả phụ thuộc vào thứ tự điều phối luồng của hệ điều
hành~\cite{goetz2006java}.\\
\textbf{Deadlock} là tình trạng hai luồng bị khóa vĩnh viễn do mỗi luồng giữ
một tài nguyên và chờ tài nguyên do luồng kia nắm giữ.
\end{dinhnghia}

Để giải quyết bài toán thống kê số lần đoán trung bình mỗi ván và trung bình
mỗi người chơi theo yêu cầu của F16, hệ thống áp dụng Java Streams
API~\cite{oraclejava21} với công thức tính toán toán học chuẩn mực:
\begin{equation}
\overline{L} = \frac{1}{|V|} \sum_{v \in V} |L_v|
\label{eq:avg}
\end{equation}
trong đó $V$ là tập hợp các ván chơi đã hoàn tất hợp lệ (loại bỏ các ván chưa
kết thúc hoặc bị hủy), và $|L_v|$ là tổng số lượt đoán ghi nhận được trong ván
chơi $v$. Thao tác tính toán~\eqref{eq:avg} được thực hiện hoàn toàn thông qua
chuỗi hàm \code{filter()}, \code{mapToInt()} và \code{average()} của
\code{IntStream} (Chương 2).

\begin{ghinho}
Trong lập trình mạng đa luồng với Java, cấu trúc \code{record} (Chương 6) bảo
đảm tính bất biến tuyệt đối (Immutability), giúp việc truyền dữ liệu giữa các
luồng mạng an toàn 100\% mà không cần sao chép tốn kém.
\end{ghinho}

\section{Công cụ và môi trường phát triển}

\begin{table}[H]
\centering
\caption{Danh mục công cụ và thư viện sử dụng trong bài tập lớn.}
\label{tab:congcu}
\begin{tabular*}{\textwidth}{@{\extracolsep{\fill}}>{\bfseries\color{daunavy}}l >{\color{daumaroon}}l p{8.4cm}}
\toprule
\rowcolor{gray100}\bfseries\color{daunavy}Thành phần & \bfseries\color{daumaroon}Phiên bản & \bfseries\color{daunavy}Vai trò trong bài tập lớn \\
\midrule
Java Development Kit   & OpenJDK 21 LTS    & Trình biên dịch và thực thi chính (hỗ trợ record, switch expression, pattern matching) \\
JUnit Platform Console & 1.10.2 Standalone & Thực thi tự động 103 ca kiểm thử TDD từ dòng lệnh~\cite{junit5docs} \\
Visual Studio Code     & 1.95+             & Môi trường lập trình (IDE), debug và quản lý workspace \\
Git / GitHub           & 2.40+             & Quản lý phiên bản mã nguồn, nhánh tính năng và worktree \\
\bottomrule
\end{tabular*}
\end{table}

\section{Các phương pháp tiếp cận liên quan}

So với phương pháp truyền đối tượng nhị phân truyền thống
(\code{ObjectInputStream} / \code{ObjectOutputStream}), giao thức văn bản
hướng dòng~\cite{fowler2002patterns} tuy đòi hỏi phân tích cú pháp chuỗi
(parsing) nhưng loại bỏ hoàn toàn nguy cơ mất an toàn bảo mật từ việc giải
tuần tự hóa (deserialization vulnerabilities) và tương thích hoàn hảo với giao
diện người dùng Swing.
